\documentclass[../../../include/open-logic-section]{subfiles}

\begin{document}

\olfileid[hi]{sth}{spine}{foundation}

\olsection{आधारिता}

हमारा काम केवल \emph{लगभग} पूरा है---\emph{पूर्णतः} समाप्त नहीं---क्योंकि
$\ZFminus$ में कुछ भी यह गारंटी नहीं देता कि \emph{प्रत्येक} समुच्चय किसी
$V_\alpha$ में है, अर्थात प्रत्येक समुच्चय किसी चरण पर बना है।

एक काफ़ी सीधा (गणितीय) अर्थ है जिसमें हमें इसकी \emph{परवाह नहीं} कि
पदानुक्रम के बाहर समुच्चय हैं या नहीं। (यदि वहाँ कोई हों, तो हम उन्हें बस
अनदेखा कर सकते हैं।) किंतु हमने समुच्चय की अपनी \emph{संकल्पना} को इस विचार
से प्रेरित किया है कि प्रत्येक समुच्चय किसी चरण पर बनता है
(\olref[z][story]{sec} में \stageshier{} देखें)। इसलिए हम पदानुक्रम के बाहर
पड़ने वाले समुच्चयों की संभावना रोकना चाहेंगे। तदनुसार हमें ऐसा नया अभिगृहीत
जोड़ना होगा जो सुनिश्चित करे कि प्रत्येक समुच्चय पदानुक्रम में कहीं आता है।

चूँकि $V_\alpha$ हमारे चरण हैं, हम बस निम्न को अभिगृहीत के रूप में जोड़ने
पर विचार कर सकते हैं:

\begin{defish}
\emph{नियमितता।} $\forall A \exists \alpha\, A \subseteq V_\alpha$
\end{defish}

यह पूरी तरह उचित दृष्टिकोण होता। किंतु अगले अनुभाग में बताए जाने वाले कारणों
से हम इसके बजाय एक वैकल्पिक अभिगृहीत अपनाएँगे:

\begin{axiom}[आधारिता]
$(\forall A \neq \emptyset)(\exists B \in A)A \cap B = \emptyset$.
\end{axiom}

कुछ प्रयास से हम ($\ZFminus$ में) दिखा सकते हैं कि आधारिता से नियमितता
निष्पन्न होती है:
\begin{defn}
प्रत्येक समुच्चय $A$ के लिए रखें:
\begin{align*}
	\text{cl}_0(A) &= A,\\
	\text{cl}_{n+1}(A) &= \bigcup \text{cl}_n(A),\\
	\text{trcl}(A) &= \bigcup_{n < \omega} \text{cl}_{n}(A).
\end{align*}
हम $\text{trcl}(A)$ को $A$ की \emph{संक्रमणशील संवृति} कहते हैं।
\end{defn}
\noindent
``संक्रमणशील संवृति'' नाम उपयुक्त है:
\begin{prop}\ollabel{subsetoftrcl}
$A \subseteq \trcl{A}$ और $\trcl{A}$ संक्रमणशील समुच्चय है।
\end{prop}

\begin{proof}
स्पष्टतः $A = \text{cl}_0(A) \subseteq \text{trcl}(A)$। और यदि
$x \in b \in \trcl{A}$, तो किसी $n$ के लिए $b \in \text{cl}_n(A)$,
अतः $x \in \text{cl}_{n+1}(A) \subseteq \text{trcl}(A)$।
\end{proof}

\begin{lem}\ollabel{lem:TransitiveWellFounded}
यदि $A$ संक्रमणशील समुच्चय है, तो कोई $\alpha$ ऐसा है कि
$A \subseteq V_\alpha$।
\end{lem}

\begin{proof}
\olref[ordinals][opps]{defsupstrict} से ``$\supstrict(X)$'' की परिभाषा याद
करते हुए दो समुच्चय परिभाषित करें:
\begin{align*}
	D  &= \Setabs{x \in A}{\forall \delta\ x \nsubseteq V_\delta}\\
	\alpha &= \supstrict\Setabs{\delta}{(\exists x \in A)
	(x \subseteq V_\delta \land (\forall \gamma \in \delta)x \nsubseteq V_\gamma)}
\end{align*}
$D = \emptyset$ मान लें। अतः यदि $x \in A$, तो कोई $\delta$ ऐसा है कि
$x \subseteq V_\delta$ और क्रमसूचकों की सुव्यवस्था से
$(\forall \gamma \in \delta)x \nsubseteq V_\gamma$; फलतः
$\delta \in \alpha$ और इसलिए
\olref[spine][Valphabasic]{Valphabasicprops} से $x \in V_\alpha$। अतः
अपेक्षानुसार $A \subseteq V_{\alpha}$।

अतः $D = \emptyset$ दिखाना पर्याप्त है। विरोधाभास द्वारा प्रमाण के लिए
अन्यथा मान लें। आधारिता से कोई $B \in D \subseteq A$ ऐसा है कि
$D \cap B = \emptyset$। यदि $x \in B$, तो $A$ की संक्रमणशीलता से
$x \in A$; और चूँकि $x \notin D$, इसलिए
$\exists \delta\ x \subseteq V_\delta$। अब रखें:
\[
	\beta = \supstrict\Setabs{\delta}{(\exists x \in b)(x \subseteq V_\delta \land (\forall \gamma < \delta)x \nsubseteq V_\gamma)}.
\]
पहले की तरह $B \subseteq V_\beta$, जो $B \in D$ दावे के विरुद्ध है।
\end{proof}

\begin{thm}\ollabel{zfentailsregularity}
नियमितता लागू होती है।
\end{thm}

\begin{proof}
$A$ को नियत करें; अब \olref{subsetoftrcl} से
$A \subseteq \trcl{A}$, जो संक्रमणशील है। अतः
\olref{lem:TransitiveWellFounded} से कोई $\alpha$ ऐसा है कि
$A \subseteq \trcl{A} \subseteq V_\alpha$।
\end{proof}

ये परिणाम दिखाते हैं कि $\ZFminus$ सशर्त
$\emph{आधारिता}\Rightarrow\emph{नियमितता}$ को सिद्ध करता है।
\olref[spine][rank]{zfminusregularityfoundation} में हम दिखाएँगे कि
$\ZFminus$, $\emph{नियमितता}\Rightarrow\emph{आधारिता}$ को सिद्ध करता है।
अतः आधारिता और नियमितता ($\ZFminus$ के सापेक्ष) \emph{समतुल्य} हैं। इसका
अर्थ है कि $\ZFminus$ को मानकर हम नियमितता से उसकी समतुल्यता देखकर आधारिता
का औचित्य दे सकते हैं। और \stageshier{} के आधार पर नियमितता का औचित्य तुरंत
दे सकते हैं।

\end{document}
